\vspace{-5pt}
\section{Theory}
We now introduce the OM action as a way to compute path probabilities under a particular stochastic differential equation (SDE), and describe its application to score-based generative modeling and transition path sampling.
\vspace{-5pt}
\subsection{Probability of paths under stochastic dynamics} 
\label{sec:stoch_dyn}
We introduce the following constant variance SDE which will underpin our proposed framework for TPS: 
%\vspace{-2ex}
\begin{equation}
    \label{eq:overdamped_langevin_main}
    \odif{\xx} = \frac{1}{\zeta}\mathbf{\Phi(x)}\odif{t} + \sqrt{2 D} \odif{\mathbf{W_t}},
\end{equation}
where $\mathbf{\Phi}(\xx) : \R^k \to \R^k$ is a drift function, $\mb{W}_t$ is a standard Weiner process, and $D, \zeta > 0$ are scalar constants governing diffusion noise levels and damping respectively. We consider drifts which can be written as the gradient of a scalar: $\mathbf{\Phi(x)} = -\frac{\partial \phi(\xx)}{\partial \xx}$, where $\phi(\xx) : \R^k \to \R$. By solving a Fokker-Planck equation for the time-varying state distribution $p(\xx, t)$, we can obtain the probability of a path $\xx(\cdot) = \left \{ \xx(t) \right \}_{t=0}^1$ sampled from the SDE in \eqref{eq:overdamped_langevin_main} (see Appendix \ref{appendix:OM_action} for complete details):
\begin{equation}
\label{eq:path_probability_main}
    P( \xx(\cdot)) \propto e^{\left(-S[ \xx(\cdot)]\right)}.
\end{equation}
To maximize this probability with respect to a path, we can equivalently minimize the negative log probability $S$, which is called the Onsager-Machlup action functional. This is the stochastic analogue of the well-known principle of least action from optics and quantum mechanics \cite{rojo2018principle}. Since we only consider discretized paths in this work, we focus on the discretized form of the OM action:
\begin{tcolorbox}[colback=cyan!15!white, boxrule=0.1mm, arc=1mm]
\begin{definition}
For a discrete path $\mathbf{X} = \{\xx^{(0)}, \dots, \xx^{(L)}\}$ generated by the SDE in \eqref{eq:overdamped_langevin_main} with drift $\mathbf{\Phi}$ and timestep size $\Delta t$, the discretized form of the \textbf{Onsager-Machlup action functional} is given by: 
\vspace{-2ex}
\begin{align}
\begin{split}
\label{eq:final_OM}
    S[\mathbf{X}] &= \frac{1}{2D} \Bigg(A[\mathbf{X}]
    + B[\mathbf{X}] +  C[\mathbf{X}] \Bigg), \\
    A[\mathbf{X}] &= \frac{1}{2\Delta t} \sum_{i=0}^{L-1} \left\|\xx^{(i+1)} - \xx^{(i)} \right\|_2^2, \\ 
    B[\mathbf{X}] &= \frac{\Delta t}{2\zeta^2} \sum_{i=1}^{L-1} \left\|\mathbf{\Phi}(\xx^{(i)})\right\|_2^2, \\ 
    C[\mathbf{X}] &= \frac{D\Delta t}{\zeta} \sum_{i=1}^{L-1} \nabla \cdot \mathbf{\Phi}(\xx^{(i)}).
\end{split}
\end{align}
\end{definition}
\end{tcolorbox}
The three summands of the OM action each have an intuitive interpretation. Term $A$ encourages smooth transitions along the discretized path. Term $B$ encourages paths remaining in regions with low-norm drifts, which are equilibria or saddle points of the underlying dynamics. Finally, term $C$ encourages paths to remain in regions with low divergence of the drift, which can be interpreted as entropically favoring regions of convexity in the landscape of $\phi$, where dynamics are more stable. The parameters $\zeta, \Delta t$, and $D$ control the relative contribution of these three terms in a physically intuitive manner. For instance, at larger values of $\Delta t$, the contribution of term $A$ is diminished, consistent with the intuition that larger ``jumps" are more probable with larger timesteps. In the limiting case of negligible diffusivity $D$ (analogous to temperature, see Appendix \ref{appendix:OM_action}), the divergence term can be omitted, yielding the Truncated OM Action:
\begin{tcolorbox}[colback=cyan!15!white, colframe=black, boxrule=0.1mm, arc=1mm]
\begin{definition}
The \textbf{truncated OM action} of a discretized path $\mathbf{X}$ is given by:
\begin{align}
\label{eq:final_OM_truncated}
S[\mathbf{X}] &= \frac{1}{2D} \big( A[\mathbf{X}] + B[\mathbf{X}] \big).
\end{align}
\end{definition}
\end{tcolorbox}
In this work, we use both the truncated and full OM actions depending on the particular system considered.  
\vspace{-3pt}
\subsection{Transition path sampling in molecular systems} 
The connection between the SDE in \eqref{eq:overdamped_langevin_main} and TPS is straightforward. Formally, in TPS we consider $d$-dimensional molecular systems with $N_p$-many particles interacting under a potential energy function $U(\mathbf{\xx}): \mathbb{R}^{N_p \times d} \rightarrow \mathbb{R}$, where $\xx \in \Omega$ is a configuration of the system, and $\Omega \in  \mathbb{R}^{N_p \times d}$ is the configuration space. The goal of TPS is to, for the given system and temperature, find most likely paths $\left\{\xx^{(i)} \right\}_{0 \leq i \leq L}$ over a time horizon $T_p$, where $\Delta t$ is the simulation timestep and $L = \frac{T_p}{\Delta t}$ is the number of discretization points in the trajectory. The trajectory traverses the two endpoints $\xx^{(0)} \in \mathcal{A} \subset \Omega$, $\xx^{(L)} \in \mathcal{B} \subset \Omega$, where $\mathcal{A}, \mathcal{B}$ typically represent distinct minima on the potential energy surface $U(\xx)$. The underlying particle dynamics are governed by the SDE in \eqref{eq:overdamped_langevin_main}, known in this context as overdamped Langevin dynamics. The scalar $\phi$ and gradient-based drift $\mathbf{\Phi}$ terms have a clear physical interpretation as the potential energy and forces, respectively:
\vspace{-5pt}
\begin{align}
    \phi(\xx) &:= U(\xx), \\
    \mathbf{\Phi}(\xx) &:= \mb{F}(\xx) = -\nabla U(\xx).
\end{align}
If the forces $\mb{F}(\xx)$ are known, the OM action (\eqref{eq:final_OM}) can be used to compute the probability of paths connecting endpoints $\xx^{(0)}, \xx^{(L)}$ under the governing dynamics.

\vspace{-3pt}
\subsection{Score-based generative modeling}
We now describe the connection between the SDE in \eqref{eq:overdamped_langevin_main} and generative models, namely denoising diffusion and flow matching. Specifically, we show that these models induce a set of stochastic dynamics whose drift is given by their learned $\textit{score}$ function. This provides a powerful framework to reason about high-probability transition paths.
\subsubsection{Stochastic dynamics under denoising diffusion models.}
\label{sec:stoch_dynamics_gen}
\emph{Denoising diffusion probabilistic models} (DDPM) \cite{ho2020denoising} are a class of score-based generative models that learn how to de-noise corrupted samples. The DDPM objective \eqref{eq:ddpm_loss} is closely linked to the score-matching objective 
\cite{vincent2011connection, song2019generative} for training a score model $s_\theta (\xx, \tau): \R^k \times \R^+ \to \R^k$ parameterized by $\theta$:
\begin{equation}
     \mathcal{L}_\text{SM}(\theta) = \mathbb{E}_{\tau, \xx \sim p_\tau} \Bigg[ \|s_\theta(\xx, \tau) - \nabla \log \left(p_\tau(\xx)\right)\|_2^2 \Bigg],
     \label{eq:score_matching}
\end{equation}
where $\nabla \log p_\tau(\xx)$ is the score of the noised marginal distribution $p_\tau$ at time $\tau$. This establishes the connection between the score and the optimal noise model with parameters $\theta^*$: $\epsilon_{\theta^*}(\xx, \tau) \propto -\nabla \log p_\tau(\xx)$. See \Cref{sec:ddpm_score_proof} for detailed statements and proofs. In order to use the OM action to compute path probabilities with a DDPM, we must construct a surrogate SDE in the form of \eqref{eq:overdamped_langevin_main}, such that paths under this SDE have high likelihood under the data distribution used to train the model. While the denoising (i.e., sampling) process of a DDPM (see \Cref{sec:ddpm_sampling}) may appear to be a natural candidate, a closer inspection reveals that it is unsuitable, as it optimizes for different likelihoods at different points of the trajectory. A large portion of the denoising trajectory thus has low likelihood under the data distribution. Therefore, we need to consider an alternative approach.
\vspace{-3pt}
\paragraph{Iterative denoising and noising as a candidate SDE.} Another hypothesis for constructing an SDE is to leverage the process of iterative one-step denoising and noising at a fixed time marginal $\tau$ of the diffusion process. Intuitively, this balances the likelihood-maximizing drift of the denoising step with the stochasticity of the noising step. Specifically, we consider the following iterated denoise-noise updates, where $\epsilon_\theta(\xx, \tau)$ is the trained denoising model from the DDPM:
\begin{align}
    \xx^{(i, \text{mid})} &= \frac{1}{\sqrt{1-\beta_\tau}}\left(\xx^{(i)} - \frac{\beta_\tau}{\sqrt{1 - \bar{\alpha}_\tau}} \epsilon_\theta(\xx^{(i)}, \tau)\right) + \sqrt{\beta_\tau} z, \\
    \xx^{(i+1)} &= \sqrt{1 - \beta_\tau} \xx^{(i, \text{mid})} + \sqrt{\beta_\tau} z',
\end{align}
where $z, z' \sim \mc N(0, I)$, $\alpha, \beta$ denote the usual diffusion model noise schedule variables, and $\bar{\alpha}_\tau = \prod_{i=1}^\tau \alpha_i$. Combining the two updates yields a single update equivalent in distribution, writing $\mb{s}_\theta(\xx, \tau) = -\left(1 / \sqrt{1-\bar{\alpha}_{\tau-1}}\right)\epsilon_\theta(\xx, \tau)$, yields:
\begin{align}
\label{eq:noise_denoise_combo}
\begin{split}
    \xx^{(i+1)} &= \xx^{(i)} + \frac{\beta_\tau\sqrt{1-\bar{\alpha}_{\tau-1}}}{\sqrt{1-\bar{\alpha}_\tau}} \mb{s}_\theta(\xx^{(i)}, \tau) + \sqrt{2\beta_\tau - \beta_\tau^2} z,
\end{split}
\end{align}
where $z \sim \mc N(0, I)$. Taking the continuum limit of DDPM sampling discretization steps to infinity, we get that $\beta_\tau \to 0$. Noting that $2\beta_\tau - \beta_\tau^2 \approx 2\beta_\tau$ and $\bar{\alpha}_\tau \to \bar{\alpha}_{\tau - 1}$ at this limit, we see that \eqref{eq:noise_denoise_combo} is an Euler-Maruyama discretization, with timestep $\beta_\tau$, of the following SDE:
\begin{equation}\label{eq:latent_sde_def}
    d{\xx} =  \mb{s}_\theta(\xx, \tau)\odif{t} + \sqrt{2} \, \odif{\mathbf{W_t}}.
\end{equation}
Note that the above construction holds for any $\tau \in \{1, \ldots, T_d\}$ where $T_d$ is the maximum diffusion time. A similar derivation can be found in \citet{arts2023two} for $\tau = 0$.

Then, \eqref{eq:latent_sde_def} is equivalent (up to constants) to \eqref{eq:overdamped_langevin_main}:
\begin{align}
    \phi(\xx) &\propto - \log p_\tau(\xx), \\
    \mathbf{\Phi}(\xx) &\propto \mb{s}_\theta(\xx, \tau) \approx \nabla \log p_\tau(\xx).
\end{align}
This means that, as introduced in \eqref{eq:path_probability_main} and \eqref{eq:final_OM}, the likelihood of discrete paths $\mathbf{X} = \left(\xx^{(t)} \right)_{0 \leq t \leq L}$ under the constructed SDE \eqref{eq:latent_sde_def} can be evaluated via the OM action defined by the trained DDPM score $\mb{s}_{\theta^*}(\xx, \tau)$:
\begin{align}
\begin{split}
    &S(\mb{X}; \theta^*) = \frac{1}{2D}\Bigg(\sum_{i=0}^{L-1}  \frac{1}{ 2\Delta t}\left\|\xx^{(i+1)} - \xx^{(i)} \right\|^2_2 \\
    &\quad 
    + \frac{\Delta t}{2\zeta^2} \left\|\mb{s}_{\theta^*}(\xx^{(i)}, \tau)\right\|_2^2 + \frac{D \Delta t}{\zeta} \ \nabla \cdot \mb{s}_{\theta*}(\xx^{(i)}, \tau)\Bigg),
\end{split}
\label{eq:generative_om_action}
\end{align}

where $\zeta = 1$, $D=1$, and $\Delta t = \beta_\tau$. Note that $\beta_\tau$ can be effectively tuned as a hyperparameter by changing the sampling discretization fidelity. In our specific scenario of TPS, we can set these parameters to physically interpretable values, with $\Delta t$, $\zeta$, and $D$ corresponding to the MD simulation timestep, friction coefficient, and diffusion coefficient, respectively (see Section \ref{sec:phys_interp}).

\subsubsection{Extension to flow matching}
We now link OM action-minimization with a broader class of generative models beyond DDPM. Flow matching models \cite{lipman2022flow} are a natural choice due to their strong performance in generative modeling tasks across modalities \cite{jing2024alphafoldmeetsflowmatching, polyak2024movie}. Similarly to DDPM, flow matching generates samples through a repeated integration process over a learned vector field. For affine flows considered in this work, the training objective for the learned velocity field $u_\theta(x, \tau)$ takes the form,
\begin{align}
    \mathcal{L}_\text{FM}(\theta) &= \mathbb{E}_{\tau, \xx \sim p_\tau} \Bigg[ \|u_\theta(\xx, \tau) - u_\tau(\xx)\|_2^2 \Bigg],\\
    u_\tau(\xx) &= \E_{\xx_0 \sim p_0, \xx_1 \sim p_1}\left[ \dot{\alpha_\tau} \xx_1 + \dot{\sigma_\tau} \xx_0 \mid \xx = \alpha_\tau \xx_1 + \sigma_\tau \xx_0\right],
\end{align}
where $p_0$ and $p_1$ are the source and target distributions, $\alpha_\tau, \sigma_\tau : [0,1] \to [0,1]$ define a curve from $\xx_0$ to $\xx_1$: $\alpha_0 = \sigma_1 = 0$, $\alpha_1 = \sigma_0 = 1$, and $\alpha_\tau, -\sigma_\tau$ are both strictly increasing functions.

While extracting the learned score $\mb{s}_\theta$ from DDPM is straightforward via $\mb{s}_\theta (\xx, \tau) = -\left(1 / \sqrt{1-\bar{\alpha}_{\tau-1}}\right)\epsilon_\theta (\xx, \tau)$, it is less clear how to do so for flow matching. However, we note the following:
\begin{enumerate}
    \item By \eqref{eq:score_matching}, the targets for the denoising model $\epsilon_\theta(\xx, \tau)$ in DDPM are equivalently the negative scores of the noised distribution, $-\nabla \log p_\tau (\xx)$.
    \item The targets $u_t(\xx)$ for the flow matching model $u_\theta(\xx, \tau)$ can be converted to scores of the flow marginal distribution $\nabla \log p_\tau (\xx)$ through the following formula (see \Cref{sec:proof_flowmatch_score} for the proof):
    \begin{equation}
        \label{eq:flow_force}
          \nabla \log p_\tau(\mathbf{x}) =  \frac{\alpha_\tau}{\dot{\sigma}_\tau \sigma_\tau \alpha_\tau - \dot{\alpha}_\tau\sigma_\tau^2} \left(\frac{\dot{\alpha}_\tau}{\alpha_\tau}\mathbf{x} - u_\tau(\mathbf{x})\right).
    \end{equation}
\end{enumerate}
\vspace{-7pt}

We can thus extract an approximate score $\mb{s}_{\theta^*}^{\mathrm{FM}}(\xx, \tau)$ from a trained flow matching model $u_{\theta^*}(\xx, \tau)$ via,
\begin{equation}
    \label{eq:flow_force_theta}
     \mb{s}_\theta^{\mathrm{FM}}(\xx, \tau) =  \frac{\alpha_\tau}{\dot{\sigma}_\tau \sigma_\tau \alpha_\tau - \dot{\alpha}_\tau\sigma_\tau^2} \left( \frac{\dot{\alpha}_\tau}{\alpha_\tau}\mathbf{x} - u_{\theta^*}(\mathbf{x}, \tau) \right).
\end{equation}

By inserting $\mb{s}_{\theta^*}^{\mathrm{FM}}(\xx, \tau)$ into the denoise-noise process defined in \eqref{eq:latent_sde_def}, we again obtain an SDE of the form of \eqref{eq:overdamped_langevin_main}. Hence, we can use the OM action to compute the log-probabilities of paths between arbitrary datapoints. 
\subsubsection{Physical interpretation for Boltzmann-distributed data}
\label{sec:phys_interp}
By assuming underlying dynamics of the form \eqref{eq:latent_sde_def}, our framework combining generative models and the OM action can be used to compute the log-probabilities of paths between samples from \textit{any} data distribution with a well-defined score. However, in the special case where the data used to train the generative model are Boltzmann distributed, i.e., $p(\mathbf{x}) \propto \exp{(-\frac{U(\mathbf{x})}{k_BT})}$, where $k_B$ is Boltzmann's constant and $T$ is the temperature, the learned score $\mb{s}_{\theta^*}(\xx, \tau = 0): \mathbb{R}^{N_p \times d} \rightarrow \mathbb{R}^{N_p \times d}$ 
is interpretable as a physical, atomistic force field:
\begin{align}
    \mb{s}_{\theta^*}(\xx, \tau =0) \approx \nabla \log p(\xx) \propto - \nabla U(\xx) = \mathbf{F}(\xx).
\end{align}
The constructed dynamics in \eqref{eq:latent_sde_def} at $\tau=0$ reduce to overdamped Langevin dynamics governing particles on a potential energy surface. Thus, we can directly set the constants in the OM action ($\Delta t$, $\zeta$, $D$, $T$) to physical values used in MD simulations for a given atomistic system. Under our OM framework, finding high-probability paths between points on a Boltzmann-distributed data manifold thus directly aligns with the conventional notion of TPS for atomistic systems.
